\documentclass[11pt]{article}
\usepackage[margin=1in]{geometry}
\usepackage{amsmath,amssymb}
\usepackage[hidelinks]{hyperref}

\title{Literature Status: Canonical AUS Injective Tensor Products}
\author{Run fa\_banach\_001, lane 16}
\date{2026-07-03}

\begin{document}
\maketitle

\section*{Status}

This is a lightweight literature-status packet.  The target question from
Garcia-Lirola and Raja, arXiv:1705.02020, is already answered in a separate
later paper by R. M. Causey, arXiv:1705.05484.  This packet is not a new proof
claim by the run.

\section*{Source Question}

Source paper: Luis Garcia-Lirola and Matias Raja, \emph{On strong asymptotic
uniform smoothness and convexity}, arXiv:1705.02020.
The relevant passage is in the introduction, page 2.  After recalling known
renorming results for injective tensor products, the paper asks whether ``the
injective tensor product of AUS spaces is an AUS space in its canonical norm.''

The same source paper proves the strong version: Corollary 4.2 says that if
\(X\) and \(Y\) are strongly AUS, then \(X \hat\otimes_\varepsilon Y\) is AUS.

\section*{Supporting Answer}

Supporting paper: R. M. Causey, \emph{Power type asymptotically uniformly
smooth and asymptotically uniformly flat norms}, arXiv:1705.05484.
The paper was submitted on 2017-05-15, after the source paper's 2017-05-04
arXiv submission.  Its abstract states the Banach-space consequence: the
injective tensor product of two AUS Banach spaces is AUS.

The decisive formal statement is Theorem 1.3.  For operators
\[
A_0:X_0\to Y_0,\qquad A_1:X_1\to Y_1,
\]
it says that if both \(A_0\) and \(A_1\) are asymptotically uniformly smooth,
then
\[
A_0\otimes A_1:
X_0\hat\otimes_\varepsilon X_1\to
Y_0\hat\otimes_\varepsilon Y_1
\]
is also asymptotically uniformly smooth.  Applying this theorem to the identity
operators \(I_X\) and \(I_Y\) gives the Banach-space statement:
\[
X,Y\text{ AUS}\quad\Longrightarrow\quad
X\hat\otimes_\varepsilon Y\text{ AUS}
\]
for the canonical injective tensor norm.

This consequence is also stated explicitly in Corollary 4.5: asymptotic uniform
smoothness passes to injective tensor products of Banach spaces.

\section*{Why This Is an Exact Literature Answer}

The supporting paper explicitly connects itself to the source paper.  In the
introduction, immediately before Theorem 1.3, Causey says that the tensor-product
result extends the recent Garcia-Lirola--Raja result on strong asymptotic
uniform smoothness; the bibliography item [13] is the source paper.

Thus the match is exact rather than merely reformulated:
\begin{itemize}
\item the source question is about the canonical injective tensor product of
plain AUS spaces;
\item Causey's theorem proves precisely that AUS passes to canonical injective
tensor products;
\item Causey explicitly knew and cited the source strong-AUS theorem as the
result being extended.
\end{itemize}

\section*{Scope Limits}

This packet does not resolve the source paper's separate Section 3 remark asking
for a reflexive Banach space with an FDD that is AUC but not strongly AUC.  It
also does not resolve the separate AUC-side question in Remark 4.3 about broader
positive classes for injective tensor products of AUC spaces.  The resolved
question here is only the canonical AUS injective tensor-product question from
the introduction.

\section*{References}

\begin{enumerate}
\item Luis Garcia-Lirola and Matias Raja, \emph{On strong asymptotic uniform
smoothness and convexity}, arXiv:1705.02020, 2017.
\item R. M. Causey, \emph{Power type asymptotically uniformly smooth and
asymptotically uniformly flat norms}, arXiv:1705.05484, 2017.
\end{enumerate}

\end{document}
